% Selected problems - Chapters 3, 4 and 5 (English)
\section*{Selected Problems}

\textbf{Chapter 3}
\begin{enumerate}
  \item[3.3] How do the surface area and volume of a sphere scale? Why? (Hint: analyze spheres of radii 1 and $R$.)
  \item[3.6] Show how the equation $y = m x^b$ becomes a linear equation on a log-log plot.
  \item[3.9] Use dimensional analysis to confirm that power equals the time rate of change of kinetic energy.
  \item[3.18] Confirm that eq. (3.28) is correct.
  \item[3.19] Show that the deflection $w_{mp}$ of a beam with bending stiffness $EI$, length $L$, under a concentrated load $P$, is governed by the two dimensionless groups in eq. (3.34).
  \item[3.28] Using reasoning similar to that which led to eq. (3.13), show that the maximum speed at which animals can run is independent of size.
  \item[3.32] It was discovered that a certain cucumber had cells that divided when they grew to 1.5 times the volume of the "remaining" cells. Cells normally divide so that the ratio between their surface and their mass remains constant. Is this a "normal" cucumber?
  \item[3.35] When the structural elements called beams vibrate... (Note: original sentence is cut/incomplete.)
\end{enumerate}

\textbf{Chapter 4}
\begin{enumerate}
  \item[4.1] Show that eq. (4.7) can be obtained by substituting $ix$ for $x$ in eq. (4.6).
  \item[4.2] Determine the first four terms of the Taylor expansions of $\tan x$ and $\cot x$ about $x=0$.
  \item[4.3] Determine the first four terms of the Taylor expansions of $\tanh x$ and $\coth x$ about $x=0$.
  \item[4.4] Use dimensional analysis to determine how the catenary parameter $c$ relates to the constant horizontal component of tension $T_0$ and the weight per unit length $\gamma$.
  \item[4.6] What does a body that weighs 10 N at Earth's surface weigh at height 10 m? At the peak of Mt. Everest?
  \item[4.10] If the gravitational potential is $V_g = -G M_E m / R$, find the exact expression as a function of altitude $z$ above Earth's surface.
  \item[4.11] Write a binomial expansion of the result of 4.10 to first order in $z$ for the potential energy above the surface.
  \item[4.12] Fill in the missing elements of the table to two-term order: $\sin x$, $\cos x$, $1-\sin x$, $1-\cos x$.
  \item[4.13] Fill in the missing elements to two-term order: $\sinh x$, $\cosh x$, $1-\sinh x$, $1-\cosh x$.
  \item[4.14] Develop a volume coefficient of expansion $\beta$ for a solid of dimensions $L_0,W_0,H_0$ paralleling the surface coefficient $\gamma$ of eq. (4.37).
  \item[4.17] Round off numbers to two significant figures as listed.
  \item[4.18] Round off numbers to three significant figures as listed.
  \item[4.19] Complete multiplications and express results to the correct significant figures.
  \item[4.20] Do 99.9 and 100.1 have the same number of significant figures? Explain.
  \item[4.21] Estimate the ranges within which each number lies: (a) 7.7 (b) 7.70 (c) 1200 (d) $1.200\times10^{-3}$.
  \item[4.22] By what percentage would the period of a pendulum change if its length was halved? Reduced by one-third? Reduced to one-third its original length?
  \item[4.24] How would the pendulum period change on Mars? On Pluto?
  \item[4.35] Estimate the error of approximating $\sin x$ with a Taylor formula to $n=4$ by evaluating the remainder $R_5$.
  \item[4.36] Do the statements $\sin x\simeq x$ and $\tan x\simeq x$ produce similar approximations? Confirm and explain.
  \item[4.37] Analog voltmeter systematic error: correct readings and compute percentage errors for given dial readings.
  \item[4.39] (a) Taylor series for $e^x$ about 0. (b) Compute $e^{0.5}$ to five significant figures using first four terms.
  \item[4.45] Analyze experimental $V$ vs $I$ data; discuss errors and eyeball a best-fit line.
  \item[4.46] Use least squares to fit $V$ versus $I$ for the same data and compare.
  \item[4.50] Far-field binomial expansion for $f(r)=\sqrt{a^2+r^2}$ when $r\gg a$.
\end{enumerate}

\textbf{Chapter 5}
\begin{enumerate}
  \item[5.8] What annual growth rate would double money in seven years?
  \item[5.9] Verify that eq. (5.12) solves the exponential differential equation given in (5.10) using $\int du/u = \ln u + C$.
  \item[5.35] Obtain actual world population numbers for 1970, 1980, 1990, 2000 and update Figures 5.1 and 5.2.
  \item[5.40] How much must be set aside in 2002 at 5.5% p.a. to reach $1{,}000{,}000$ by 2022? By 2042?
  \item[5.47] Iwo Jima: with $F_0=54000$, $E_0=21500$, $\lambda_F=0.0106$, $\lambda_E=0.0544$ per day — how long and how many survivors for the victor?
\end{enumerate}

% End of file (English)
\begin{enumerate}
    \item[\textbf{Problem 3.1.}] On what basis did the Lilliputians conclude that Gulliver needed 1728 times as much food as they did?

    \item[\textbf{Problem 3.2.}] How would the Lilliputians' conclusion change if they had thought about the exchange of energy between a person (of any size) and the surrounding environment?

    \item[\textbf{Problem 3.3.}] How do the surface area and volume of a sphere scale? Why? (Hint: Analyze spheres of radii 1 and R.)

    \item[\textbf{Problem 3.4.}] Explain what would happen to an angle between two lines inscribed on a balloon as it was inflated to a radius R from a radius of 1.

    \item[\textbf{Problem 3.5.}] Confirm that eq. (3.2) does adequately portray the straight line drawn in Figure 3.2.

    \item[\textbf{Problem 3.6.}] Show how the equation $y = mx^b$ becomes a linear equation in a log-log plot.

    \item[\textbf{Problem 3.7.}] Write eq. (3.6) in a form suitable for a log-log plot.

    \item[\textbf{Problem 3.8.}] Confirm the dimensional relationship of eq. (3.10).

    \item[\textbf{Problem 3.9.}] Use dimensional analysis to confirm that power is equal to the time rate of change of kinetic energy.

    \item[\textbf{Problem 3.10.}] Confirm the dimensional relationship of eq. (3.12).

    \item[\textbf{Problem 3.11.}] Confirm that eq. (3.13) does show that $v$ is independent of $L$.

    \item[\textbf{Problem 3.12.}] What is the hearing range of an elephant? A whale? How do these ranges compare with those of humans?

    \item[\textbf{Problem 3.13.}] Confirm that the dimensions of eq. (3.15) are 1/T.

    \item[\textbf{Problem 3.14.}] Confirm that the dimensions of eq. (3.16) are 1/T.

    \item[\textbf{Problem 3.15.}] Confirm that the dimensions of eq. (3.17) are 1/T.

    \item[\textbf{Problem 3.16.}] Under what conditions is eq. (3.24) dimensionally consistent?

    \item[\textbf{Problem 3.17.}] Confirm that the stress of eq. (3.27) satisfies eq. (3.24).

    \item[\textbf{Problem 3.18.}] Confirm that eq. (3.28) is correct.

    \item[\textbf{Problem 3.19.}] Show that the deflection $w_{mp}$ of a beam with bending stiffness $EI$, length $L$, and under a concentrated load $P$ is governed by the two dimensionless groups in eq. (3.34).

    \item[\textbf{Problem 3.20.}] Why is the torque, $\tau$, in the apparatus of Figure 3.12 a constant?

    \item[\textbf{Problem 3.21.}] Expressed in terms of the wheel's geometric and gravitational properties, what is the magnitude of the torque in Problem 3.20?

    \item[\textbf{Problem 3.22.}] Confirm that eq. (3.30) is dimensionally correct.

    \item[\textbf{Problem 3.23.}] Confirm that eq. (3.31) is dimensionally correct.

    \item[\textbf{Problem 3.24.}] Calculate and confirm the estimated number of revolutions in the last column of Table 3.1.

    \item[\textbf{Problem 3.25.}] Confirm that eq. (3.39) is the correct representation of eq. (3.37) in terms of the five scaling factors of a simple beam.

    \item[\textbf{Problem 3.26.}] Formulate a hypothesis to explain why a wood pigeon and a buzzard appear to have such different proportions of $W_{fm} / W_b$ in Figure 3.2.

    \item[\textbf{Problem 3.27.}] Show that the equation that describes the log-log plot of Figure 3.7 can be found to be $h \simeq 1.23l^{0.68}$, where $h$ and $l$ are, respectively, the nave height and church length rendered dimensionless by dividing each by 1 ft.

    \item[\textbf{Problem 3.28.}] Using reasoning similar to that which brought us to eq. (3.13), show that the maximum speed at which animals can run is independent of size.

    \item[\textbf{Problem 3.29.}] The velocity of blood in the aorta is related to the difference in pressure between the heart and the arteries. Find the relationship between the velocity of the blood and the pressure difference. (Hint: Use the work-energy theorem as we did for bird hovering in Section 3.3.2.)

    \item[\textbf{Problem 3.30.}] The stilt, a small long-legged bird, was described in \textit{Gulliver's Travels} as weighing 4.5 ounces and having legs 8 inches long. A flamingo has a similar shape and weighs 4 pounds. Apply scaling arguments to show that its legs should be about 20 inches long (as they actually are!).

    \item[\textbf{Problem 3.31.}] Given that a robin weighs about 2 ounces, could we scale the length of its legs from the stilt data provided in Problem 3.30? Explain your answer.

    \item[\textbf{Problem 3.32.}] It was discovered that a certain cucumber had cells that divided when they grew to 1.5 times the volume of the "remaining" cells. Cells normally divide so that the ratio between their surface and their mass remains constant. Is the cucumber described as a "normal" cucumber?

    \item[\textbf{Problem 3.33.}] Find the range of values of the variable $x$ for which the approximation $\cosh(x/\lambda) \simeq \frac{1}{2}e^{x/\lambda}$ is acceptable, for scaling factors $\lambda = 1$ and $\lambda = 6$. Plot both functions for each of the two scaling factors.

    \item[\textbf{Problem 3.34.}] An experiment to determine the natural or fundamental period of oscillation of a simple spring-mass system (see Figure P3.34) is set up as follows. A spring of stiffness $k$ is fixed at one end and connected to a mass $m$ at its other, with the mass being able to move along an ideal, frictionless air track. The mass is displaced a distance $x_0$ from its initial resting position, after which it oscillates along the air track around that initial position. The time needed for a complete oscillation-that is, the period-is measured several times for several periods in succession, with the results being compared to the theoretical formula for the period, $T$: $T = 2\pi\sqrt{m/k}$. Assuming that $k$ and $m$ are known and that the stopwatch used to measure the period is accurate to $\pm 1\%$. What are the possible pitfalls that might prevent the successful experimental determination of $T$?

    \item[\textbf{Problem 3.35.}] When the structural elements called beams vibrate... \textit{(Note: This sentence is cut off/incomplete in the original text).}

    \item[\textbf{Problem 3.36.}] A steel beam of length of 20 cm is to be used to model a prototype timber beam whose span is 3.6 m. (a) Verify that the dimensionless group containing the load, the modulus, and the length is $P/EL^2$. (b) If the timber beam is to carry a load of 9000 N at a point 1.5 m from the left end, what load must be applied to the model to determine whether the prototype can carry its intended load? (Assume that the load-carrying capacity is the only behavior of interest here.)

\end{enumerate}
